\SourceLeaf{147}%
zustand her. Zu der Zeit also, wo die Witwenzahl in der idealen Gesellschaft das Maximum 1300 ganz oder fast ganz erreicht haben würde, wird in der wirklichen ein Auf- und Abschwanken über und unter diese Zahl hinaus, Statt finden. Es werden zu einer Zeit 1350, auch wohl 1360, da sein können, zu einer andern 1250, auch wohl nur 1240. Allein ganz unmöglich ist es, hier scharfe Grenzen zu setzen, und von irgend einer Zahl, man wähle welche man wolle, mit Bestimmtheit zu behaupten, dass diese zwar noch erreicht, die nächst höhere aber nicht mehr erreicht werden könne. Die Wahrscheinlichkeitsrechnung, in ihrem heutigen Zustande, lehrt solche Schwankungen auf ein bestimmtes Maass zurückzuführen, welches aber nicht von dem Extrem hergenommen wird, sondern von der Berücksichtigung aller Zwischenstufen und ihrer relativen Wahrscheinlichkeit. Natürlich ist hier durchaus nicht der Ort, diese Theorie weiter zu verfolgen; auch hängt ihre Anwendbarkeit auf concrete Fälle davon ab, dass man entweder eine mathematisch präcise Kenntniss von den bedingenden Elementen habe (wie bei Glücksspielen z. B. gewöhnlich der Fall ist), oder dass ein zureichender Reichthum von Erfahrungen zu Gebote stehe. Beides trifft aber in Beziehung auf solche Gegenstände, wie Witwenkassen sind, nicht zu, so dass a priori für eine solche Bemessung nichts geschehen kann. Aber das lehrt doch die Theorie: aus den Schwankungen, die in einem Falle vorgekommen sind, auf diejenigen zu schliessen, die mit gleichem Recht in einem andern qualitativ gleichen aber quantitativ verschiedenen Falle erwartet werden müssen. Nach ihrer absoluten Grösse werden in einer grossen Gesellschaft die Schwankungen grösser sein, als in einer kleinen; nach der relativen aber wird es sich umgekehrt verhalten. Wenn also z. B. eine Gesellschaft von obigem Zuschnitt zu einer Zeit bis 60 Witwen mehr, zu einer andern bis 60 weniger Witwen gezählt hätte, als die Durchschnittsgrösse 1300, so jedoch, dass die Zahl 1360, oder eine ihr nahe kommende, mit einiger Andauer vorgekommen wäre, und eben so die Zahl 1240 oder eine ihr nahe liegende, so würde man schliessen dürfen, dass in einer andern Gesellschaft, die bei sonst ähnlichen Verhältnissen durchschnittlich nur 26 Ehepaare zählt, die Normalwitwenzahl eben so leicht um 6 vermehrt oder vermindert erscheinen könne. Oder, etwas anders ausgedrückt: Eben so gut, wie in der grossen Gesellschaft zwischen 1360 unn 1240, wird die Witwenzahl in der kleinen zwischen 7 und 19 auf und abschwanken können; das \textit{absolute} Schwanken ist, wenn der Umfang der grossen Gesellschaft 100 mal grösser ist wie der Umfang der kleinen, in der grossen zehnmal so gross wie in der kleinen; mit dem \textit{relativen} Schwanken, welches in der grossen Gesellschaft (nach obigen Zahlen) \(4\tfrac8{13}\) Procent betragen würde und bei der kleinen \(46\tfrac2{13}\) Procent, verhält es sich gerade umgekehrt.

So viel von der Sache. Was den Namen betrifft, so habe ich, eben, für die Zahlen der Beispiele 1300 und 13, anstatt der weitschweifigen Umschreibung ‘Durchschnittszahl der Witwen nach eingetretenem Beharrungszustande’ die Benennung ‘Normalwitwenzahl’ gebraucht, welche man leicht als nicht unpassend gelten lassen wird: aber \textit{üblich} ist sie meines Wissens nicht. Es ist vielmehr ganz gewöhnlich, jenen Begriff kurzweg mit \textit{höchster Witwenzahl} zu bezeichnen, wobei von den regellosen Fluctuationen ganz abstrahirt, und nur der Beharrungszustand im Allgemeinen im Gegensatz zu der vorhergegangenen Zeit berücksichtigt wird. Für einen einigermassen Sachverständigen wird hiebei ein Missverständniss nicht leicht möglich sein; jedenfalls aber ist so viel gewiss, und aus dem Gutachten, wenn man es nur mit einiger Aufmerksamkeit liest, sogleich zu erkennen, dass \textsc{Kritter} in demselben die höchste Witwenzahl in \textit{diesem} Sinn und nur in diesem Sinn verstanden hat. P. hingegen dachte sich dabei etwas ganz anderes, nemlich die Zahl, über welche die Anzahl der Witwen nach höchster Wahrscheinlichkeit niemals sollte hinausgehen können. Beide Zahlen vermengen, ist ungefähr dasselbe, als wenn man den Durchschnittspreis eines Handelsartikels mit dem höchsten Preise verwechselte. Es ist hiedurch die oben S. [140] unter 3 gemachte Ausstellung hinreichend bewiesen. Was aber eine Frage nach der höchsten Witwenzahl im P.’schen Sinne des Worts betrifft, so ist sie eine solche, auf welche eine bestimmte Antwort sich gar nicht geben lässt.
